% Original Girard1987 route; no appeal to modern switching sequentialization.
\section{Girard's trip intervals, empires, and maximal splitting}
\label{sec:p07-girard-trips-empires}

The hypotheses in this section are those of Girard's cut-free
multiplicative proof structures: finitely many separately identified
occurrences, general-formula dual axiom links, tensor and par links,
one producer and at most one consumer per occurrence, and no units,
boxes, cuts, open hypotheses or Mix. At every setting of \emph{both}
tensor and par switches the successor permutation has one long orbit.
We use the original port clauses on \cref{sheet:p07-mll-trips}, not
par-switching acyclicity. The proof follows the original trip-and-empire
dependency \cite[pp.~35--40]{Girard1987AuthorHosted}.

\begin{hbdefinition}{def:p07-girard-empire-interval}{Cyclic intervals and empire gates}
Write \(I_\sigma(X)\) for the closed forward interval from \(X^+\)
to \(X^-\) in the single trip at setting \(\sigma\).
If \(X\) is a premise of a fixed tensor, its empire is
\[
 eX=\{F:\{F^+,F^-\}\subseteq I_\sigma(X)
                      \text{ for every setting }\sigma\}.
\]
For a set \(E\) of occurrences write
\(\widehat E=\{F^+,F^-:F\in E\}\).
These are sets of occurrences and ports, respectively; an interval
containing one gate does not put that occurrence in an empire.
The notation \(eX\) always refers to the unique tensor consuming the
specified occurrence \(X\), not to a formula label used elsewhere.
\end{hbdefinition}

\begin{hblemma}{thm:p07-girard-empire-local-blocks}
 {proved=annex:p07.girard-empire.local-blocks}{The full local switch-change calculation}
Fix all settings except that of a link with premises \(C,D\) and root \(R\).
Deleting its three switch-dependent successor arcs leaves three disjoint
directed paths covering all gates. They are unchanged by flipping that link.
For tensor the paths have endpoints
\[
 [C^+,C^-],\qquad[D^+,D^-],\qquad[R^-,R^+].
\]
The \(L\) setting orders them as \(C,R,D\); the \(R\) setting as \(C,D,R\).
For par their endpoints are
\[
 [C^+,D^-],\qquad[D^+,C^-],\qquad[R^-,R^+],
\]
and the settings order them as \(C,D,R\) and \(C,R,D\), respectively.
All six possible endpoint matchings are accounted for in the proof.
These normal forms include the passage-order conclusions of the
original Lemmas~2.9.1--2.9.2.
\end{hblemma}

\begin{hblemma}{thm:p07-girard-empire-interval-cases}
 {proved=annex:p07.girard-empire.interval-cases}{The interval positions in both block orders}
Consider three disjoint nonempty paths in the two cyclic orders
\(P,Q,R\) and \(P,R,Q\), with unchanged internal order.
Let \(a^+,a^-\) be distinct non-endpoint gates of these paths.
If both endpoints of one specified path belong to
\([a^+,a^-]\) in \emph{both} orders, then \(a^+,a^-\) lie in the
same path with \(a^-\) earlier than \(a^+\) internally.
Consequently the interval contains all six path endpoints in both orders.
The proof covers both same-path orientations and every distinct-path
placement; it assumes no empire closure.
\end{hblemma}

\begin{hblemma}{thm:p07-girard-empire-closure}
 {proved=annex:p07.girard-empire.closure}{All seven original empire closure facts}
For the fixed tensor \(A,B\) with root \(T=A\otimes B\):
\begin{enumerate}
\item \(A\in eA\).
\item \(eA\cap eB=\varnothing\), and \(T\notin eA\cup eB\).
\item If \(C\in eA\) and an axiom joins \(C\) with its dual occurrence
      \(C'\), then \(C'\in eA\).
\item If \(R=C\otimes D\in eA\), then \(C,D\in eA\).
\item If \(R=C\bindnasrepma D\in eA\), then \(C,D\in eA\).
\item If \(R=C\otimes D\) is produced by a tensor \emph{different}
      from the fixed \(T\)-link, and \(C\in eA\) or \(D\in eA\),
      then \(R\in eA\).
\item If \(C,D\in eA\) are the premises of a par with root
      \(R=C\bindnasrepma D\), then \(R\in eA\).
\end{enumerate}
The same assertions hold with \(A,B\) exchanged.
The distinguished-tensor exception and the cases \(R=A\) are
explicitly included, rather than suppressed in the switch argument.
\end{hblemma}

\begin{hbtheorem}{thm:p07-girard-empire-realization}
 {proved=annex:p07.girard-empire.realization}{Trip Theorem 2.9.5, with its free settings}
There is a setting \(\sigma\) with
\(I_\sigma(A)=\widehat{eA}\).
At a par with exactly one premise in \(eA\), set the switch so that
that premise reflects from its minus gate to its own plus gate:
right if the inside premise is left, and left if it is right.
Every other switch may be prescribed arbitrarily before completing
the setting. In particular a producer path wholly within \(eA\)
can be selected without violating the boundary prescriptions.
\end{hbtheorem}

\begin{hbdefinition}{def:p07-girard-empire-ancestry}{Occurrence ancestry, not syntactic substrings}
An occurrence is upstream of another when a backwards chain of
producing-connective premise incidences reaches it.
We also explicitly allow a zero-length chain when stated.
The syntactic subformulas of an opaque general-formula axiom label
are not additional occurrences or hereditary premises.
\end{hbdefinition}

\begin{hbcorollary}{thm:p07-girard-empire-ancestor}
 {proved=annex:p07.girard-empire.ancestor}{Corollary 2.9.6 and the zero-length boundary}
If \(R=C\otimes D\) is \(A\) itself or is upstream of \(A\), then
\(eC\cup eD\subseteq eA\). The zero-length case is stated explicitly
because a tensor can feed a terminal tensor directly.
\end{hbcorollary}

\begin{hblemma}{thm:p07-girard-empire-joint}
 {proved=annex:p07.girard-empire.joint}{The simultaneous-realization obstruction}
Let \(T=A\otimes B\) be terminal. If there is no par with one
premise in \(eA\) and the other in \(eB\), both empires can be
realized at the same setting. It follows that
\[
 N=eA\mathbin{\dot\cup}eB\mathbin{\dot\cup}\{T\}
\]
as occurrence sets. Therefore failure of this covering forces
an actual crossing par; it is not inferred from maximality alone.
\end{hblemma}

\begin{hblemma}{thm:p07-girard-empire-crossing}
 {proved=annex:p07.girard-empire.crossing}{Both crossing-par switch cases yield a larger empire}
Suppose a par \(S\) has one premise in \(eA\) and the other in
\(eB\), in either ordered position. Call the respective occurrences
\(C,D\), retaining their actual left/right positions. Let \(J\) be a premise of a terminal tensor \(U\)
such that \(S=J\) or \(S\) is upstream of \(J\).
Then
\[
 eA\cup eB\subsetneq eJ\subseteq
       eH\cup eJ,
\]
where \(H\) is the other premise of \(U\); its actual formula is
\(H\otimes J\) or \(J\otimes H\) according to the retained port order.
The first inclusion is proved using two separately realized settings:
one selects \(C\) at \(S\), the other selects \(D\).
No single-setting nesting is assumed to establish the universal empires.
\end{hblemma}

\begin{hbtheorem}{thm:p07-girard-empire-splitting}
 {proved=annex:p07.girard-empire.splitting}{Original maximal-empire Splitting Theorem 2.9.7}
If a trip-correct structure has a terminal tensor and no terminal par,
there exists a terminal tensor \(T=A\otimes B\) with
\[
 N=eA\mathbin{\dot\cup}eB\mathbin{\dot\cup}\{T\}.
\]
Choose a tensor whose \(eA\cup eB\) is maximal under inclusion among
\emph{all} tensor links. The proof establishes its terminality and
covering from the preceding lemmas, with no splitting premise.
\end{hbtheorem}

\begin{hblemma}{thm:p07-girard-empire-restriction}
 {proved=annex:p07.girard-empire.restriction}{The two restricted structures are themselves original nets}
At the splitting tensor, restricting to \(eA\) and to \(eB\)
after removing \(T\) produces two nonempty original trip-correct
proof structures. Their respective distinguished new terminal
occurrences are \(A,B\). Producer closure, the absence of a remaining
crossing link and the correctness of \emph{every} restricted setting
are part of the proof.
\end{hblemma}

\begin{hbtheorem}{thm:p07-girard-empire-sequentialization}
 {proved=annex:p07.girard-empire.sequentialization}{The full original cut-free sequentialization argument}
Every original cut-free multiplicative long-trip net is the occurrence
construction of a closed Mix-free sequent proof with general-formula
axioms. Terminal-par removal and maximal-empire tensor splitting
supply all induction cases. This converse uses no modern switching
criterion or modern internal-par splitter.
\end{hbtheorem}

\begin{proofboundary}{def:p07-girard-empire-scope}{What the original empire completion covers}
These proofs reconstruct the original unit-free cut-free trip route,
including the restriction step omitted from the printed proof.
They do not insert Mix, units, proof boxes or a modern correctness
criterion into the hypotheses. Original CUT-marker trips and their
preservation, the generic cyclic-door box extension, all PN1/PN2
contractions and the separately identified 1989 presentation retain
their own full-source duties \cite[pp.~40--47]{Girard1987AuthorHosted}.
\end{proofboundary}
\index{proof net!original trip criterion}
\index{empire!maximal splitting}
